\section{FedAvg proof and matched attribution}
\label{app:fa-proof}

For FedAvg, persistent controls are absent, so the common estimates of
Appendix~\ref{app:scaf-proof} apply with $\xi_c=0$. We verify their
inputs and then compute the local second moment that determines the FedAvg
coefficient. All expectations are under the FedAvg invariant law, with
$s,r,v$ defined in Equation~\eqref{eq:scales}.

\subsection{Application of the common estimates}

The stationary moment estimates of \citet[Proposition~3 and
Lemmas~5--6]{mangold2025fedavg} give
Equation~\eqref{eq:A-coarse-even} under the step-size threshold in
Appendix~\ref{app:common-inputs}. Equation~\eqref{eq:fa-local} is the
common local recursion with $\xi_c=0$, and the independent fresh noise
satisfies the assumptions of the common lemmas. The two round identities
used there become
\begin{align}
 x^+-x&=-\gamma\sum_{h=0}^{H-1}\frac1N\sum_c\nabla f(\theta_{c,h})
         -\gamma\sum_{j=1}^H\bar\eps_j,
 \label{eq:B-global-increment}\\
 x^+&=P^Hx-\gamma\sum_{j=1}^HP^{H-j}\bar\eps_j
             -\gamma\sum_{h=0}^{H-1}P^{H-1-h}\bar q_h.
 \label{eq:B-round-exact}
\end{align}
These identities establish the same stationary balance and nonlinear
round decomposition used in Appendix~\ref{app:scaf-proof}. Applying
Lemmas~\ref{lem:common-coarse} and~\ref{lem:common-higher} first gives
the required local and higher-moment bounds. Lemma~\ref{lem:A-sharp-second-moment}
then identifies the global raw second moment in Equation~\eqref{eq:A-sharp-Mx},
and Lemma~\ref{lem:common-taylor} supplies the stationary mean-balance
identity in Equation~\eqref{eq:A-master}. Only the local second-moment
increment remains to be evaluated.

\subsection{Local coefficient and completion of the proof}

Without persistent controls, $Q_\xi=0$ and the stochastic part of the
local displacement contains only fresh noise. Equations~\eqref{eq:A-DR}
and~\eqref{eq:A-DD} therefore reduce to
\[
 D_{c,h}=-\gamma\sum_{j=1}^h\eps_{c,j},\qquad
 \frac1N\sum_c\E[xD_{c,h}^\top]=0,\qquad
 \frac1N\sum_c\E[D_{c,h}D_{c,h}^\top]=\gamma^2hC.
\]
The already proved remainder estimate in Lemma~\ref{lem:common-coarse} yields
\begin{equation}
\label{eq:B-U}
 U_h=\gamma^2hC+O(r),\qquad
 \sum_{h=0}^{H-1}U_h=\gamma^2\frac{H(H-1)}2C+O(r).
\end{equation}
Substituting this and $M_x=(\gamma/N)S_C+O(r)$ into
Equation~\eqref{eq:A-master} gives
\begin{equation}
\label{eq:B-final-expansion}
 b^{\FedAvg}_{\gamma,N,H}
 =-\frac{\gamma}{2N}A^{-1}T[S_C]
 -\gamma^2\frac{H-1}{4}A^{-1}T[C]
 +O_{H,d}\!\left(\frac{\gamma^2}{N}+\gamma^3\right),
\end{equation}
proving Theorem~\ref{thm:fa-main}.
Subtracting the two explicit second-order coefficients gives
\[
 B_{20}^{\Scaf}(H)-B_{20}^{\FedAvg}(H)
 =-\frac{(H-1)(2H-1)}{12H}A^{-1}T[C]
 =B_{20}^{\mathrm{ctrl}}(H).
\]
The common $\gamma/N$ term cancels in the matched comparison, while
the difference of the two uniform remainders remains
$O_{H,d}(\gamma^2/N+\gamma^3)$. This completes
Corollary~\ref{cor:attribution}.
